\documentclass[12pt,a4paper,twocolumn]{ctexart}
\usepackage[landscape]{geometry}
\geometry{top=1.5cm,bottom=1.5cm,left=1.5cm,right=1.5cm}
\usepackage{amsmath,amssymb}
\usepackage{array}
\usepackage{multicol}

\setlength{\columnsep}{1cm}
\pagestyle{empty}
\title{\textbf{数学 · 高考速查速记}}
\date{}
\begin{document}
\maketitle

%----------- 集合 -----------
\subsection*{集合}
$A\cap B=\{x\mid x\in A\text{且}x\in B\}$；
$A\cup B=\{x\mid x\in A\text{或}x\in B\}$ \\
$\complement_U(A\cap B)=\complement_U A\cup\complement_U B$ \\
$n$元集子集$2^n$个，真子集$2^n-1$个

%----------- 复数 -----------
\subsection*{复数}
$z=a+bi,\;i^2=-1,\;\bar{z}=a-bi,\;|z|=\sqrt{a^2+b^2}$ \\
$(a+bi)\pm(c+di)=(a\pm c)+(b\pm d)i$ \\
$(a+bi)(c+di)=(ac-bd)+(ad+bc)i$ \\
$\dfrac{a+bi}{c+di}=\dfrac{(a+bi)(c-di)}{c^2+d^2}$
\begin{small}
$i$周期: $i^1=i,i^2=-1,i^3=-i,i^4=1$；$z\bar{z}=|z|^2$
\end{small}

%----------- 不等式 -----------
\subsection*{不等式}
$\dfrac{a+b}{2}\ge\sqrt{ab}\;(a,b>0)$ 当$a=b$取等\\
柯西: $(a^2+b^2)(c^2+d^2)\ge(ac+bd)^2$ \\
绝对值: $|a|-|b|\le|a\pm b|\le|a|+|b|$

%----------- 函数 -----------
\subsection*{函数性质}
奇: $f(-x)=-f(x)$；偶: $f(-x)=f(x)$ \\
$f(x+a)=f(x)\Rightarrow T=a$ \\
$f(x+a)=-f(x)\Rightarrow T=2a$ \\
$f(x+a)=\frac1{f(x)}\Rightarrow T=2a$

%----------- 导数 -----------
\subsection*{导数公式}
\begin{tabular}{ll}
$(x^n)'=nx^{n-1}$ & $(\sin x)'=\cos x$ \\
$(\cos x)'=-\sin x$ & $(\tan x)'=\sec^2x$ \\
$(e^x)'=e^x$ & $(a^x)'=a^x\ln a$ \\
$(\ln x)'=\frac1x$ & $(\log_ax)'=\frac1{x\ln a}$
\end{tabular}

\textbf{运算法则}: $(uv)'=u'v+uv'$；$\big(\frac{u}{v}\big)'=\frac{u'v-uv'}{v^2}$；$[f(g(x))]'=f'(g(x))g'(x)$

\textbf{切线放缩}: $e^x\ge x+1$；$\ln x\le x-1$；$\sin x\le x\;(x\ge0)$

%----------- 三角 -----------
\subsection*{三角恒等变换}
$\sin^2\alpha+\cos^2\alpha=1$；$\tan\alpha=\dfrac{\sin\alpha}{\cos\alpha}$ \\
$\sin(\alpha\pm\beta)=\sin\alpha\cos\beta\pm\cos\alpha\sin\beta$ \\
$\cos(\alpha\pm\beta)=\cos\alpha\cos\beta\mp\sin\alpha\sin\beta$ \\
$\sin2\alpha=2\sin\alpha\cos\alpha$ \\
$\cos2\alpha=\cos^2\alpha-\sin^2\alpha=2\cos^2\alpha-1=1-2\sin^2\alpha$ \\
$a\sin x+b\cos x=\sqrt{a^2+b^2}\sin(x+\varphi),\;\tan\varphi=\frac{b}{a}$

\subsection*{解三角形}
$\frac{a}{\sin A}=\frac{b}{\sin B}=\frac{c}{\sin C}=2R$ \\
$a^2=b^2+c^2-2bc\cos A$ \\
$S=\frac12bc\sin A=\frac{abc}{4R}$

%----------- 数列 -----------
\subsection*{数列}
等差: $a_n=a_1+(n-1)d$；$S_n=\frac{n(a_1+a_n)}{2}=na_1+\frac{n(n-1)}{2}d$ \\
等比: $a_n=a_1q^{n-1}$；$S_n=\frac{a_1(1-q^n)}{1-q}\;(q\neq1)$

裂项: $\frac1{n(n+1)}=\frac1n-\frac1{n+1}$；$\frac1{(2n-1)(2n+1)}=\frac12(\frac1{2n-1}-\frac1{2n+1})$

%----------- 向量 -----------
\subsection*{平面向量}
$\vec{a}\cdot\vec{b}=x_1x_2+y_1y_2=|\vec{a}||\vec{b}|\cos\theta$ \\
$\vec{a}\perp\vec{b}\iff\vec{a}\cdot\vec{b}=0$；$\vec{a}\parallel\vec{b}\iff x_1y_2=x_2y_1$ \\
$|\vec{a}|=\sqrt{x^2+y^2}$；投影: $\frac{\vec{a}\cdot\vec{b}}{|\vec{b}|}$

%----------- 立体几何 -----------
\subsection*{立体几何}
线面角: $\sin\theta=\frac{|\vec{v}\cdot\vec{n}|}{|\vec{v}||\vec{n}|}$ \\
二面角: $\cos\theta=\frac{|\vec{n}_1\cdot\vec{n}_2|}{|\vec{n}_1||\vec{n}_2|}$（锐角取正）\\
点面距: $d=\frac{|\overrightarrow{AP}\cdot\vec{n}|}{|\vec{n}|}$ \\
体积: $V_{\text{柱}}=Sh,\;V_{\text{锥}}=\frac13Sh,\;V_{\text{球}}=\frac43\pi R^3$

%----------- 解析几何 -----------
\subsection*{解析几何}
椭圆: $\frac{x^2}{a^2}+\frac{y^2}{b^2}=1$，$c^2=a^2-b^2$，$e=\frac{c}{a}\in(0,1)$ \\
双曲线: $\frac{x^2}{a^2}-\frac{y^2}{b^2}=1$，$c^2=a^2+b^2$，$e=\frac{c}{a}>1$，渐近线$y=\pm\frac{b}{a}x$ \\
抛物线: $y^2=2px$，焦点$(\frac{p}{2},0)$，准线$x=-\frac{p}{2}$，$e=1$ \\
弦长: $|AB|=\sqrt{1+k^2}|x_1-x_2|$

%----------- 概率统计 -----------
\subsection*{概率统计}
$P(B|A)=\frac{P(AB)}{P(A)}$；$P(B)=\sum P(A_i)P(B|A_i)$ \\
二项$B(n,p)$: $P(X=k)=C_n^kp^k(1-p)^{n-k}$，$E(X)=np$，$D(X)=np(1-p)$ \\
正态$N(\mu,\sigma^2)$: $P(\mu-\sigma<X<\mu+\sigma)\approx0.6827$

%----------- 排列组合 -----------
\subsection*{排列组合}
$A_n^m=\frac{n!}{(n-m)!}$；$C_n^m=\frac{n!}{m!(n-m)!}$ \\
$C_n^m=C_n^{n-m}$；$C_n^m+C_n^{m+1}=C_{n+1}^{m+1}$ \\
$(a+b)^n=\sum_{k=0}^n C_n^ka^{n-k}b^k$；$T_{k+1}=C_n^ka^{n-k}b^k$

%----------- 数值速查 -----------
\subsection*{常用数值}
$\sqrt2\approx1.414,\;\sqrt3\approx1.732,\;\sqrt5\approx2.236$ \\
$\pi\approx3.142,\;e\approx2.718$ \\
$\ln2\approx0.693,\;\ln3\approx1.099$ \\
$\sin\frac{\pi}{6}=\frac12,\;\sin\frac{\pi}{4}=\frac{\sqrt2}2,\;\sin\frac{\pi}{3}=\frac{\sqrt3}2$ \\
$\cos\frac{\pi}{6}=\frac{\sqrt3}2,\;\cos\frac{\pi}{4}=\frac{\sqrt2}2,\;\cos\frac{\pi}{3}=\frac12$

\end{document}
